\chapter{Linear Algebra for Quantum States}

\section{Vectors as lists and arrows}

A vector in \(\R^n\) is a list
\[
  \bm v=(v_1,\ldots,v_n).
\]
The dot product is
\[
  \bm u\cdot\bm v=\sum_{i=1}^n u_iv_i.
\]
The length is \(\norm{\bm v}=\sqrt{\bm v\cdot\bm v}\).  The Cauchy--Schwarz inequality,
\[
  |\bm u\cdot\bm v|\le \norm{\bm u}\norm{\bm v},
\]
means that the angle formula
\[
  \bm u\cdot\bm v=\norm{\bm u}\norm{\bm v}\cos\theta
\]
is sensible.  One proof begins with the nonnegative quadratic
\[
  0\le \norm{\bm u-\lambda \bm v}^2
  =
  \norm{\bm u}^2-2\lambda\,\bm u\cdot\bm v+\lambda^2\norm{\bm v}^2.
\]
Its discriminant cannot be positive, so
\[
  4(\bm u\cdot\bm v)^2-4\norm{\bm u}^2\norm{\bm v}^2\le0.
\]

\section{Matrices as linear machines}

A matrix \(A\) maps vectors to vectors:
\[
  (A\bm v)_i=\sum_j A_{ij}v_j.
\]
Linearity means
\[
  A(\alpha\bm u+\beta\bm v)=\alpha A\bm u+\beta A\bm v.
\]
Matrix multiplication is composition:
\[
  (AB)_{ij}=\sum_k A_{ik}B_{kj}.
\]
Usually \(AB\ne BA\).  Their difference,
\[
  [A,B]=AB-BA,
\]
is the commutator.  Quantum mechanics is full of commutators because the order of operations matters.

\section{Eigenvectors}

An eigenvector is a nonzero vector whose direction is unchanged:
\[
  A\bm v=\lambda\bm v.
\]
The number \(\lambda\) is an eigenvalue.  This equation is equivalent to
\[
  (A-\lambda I)\bm v=0.
\]
A nonzero solution exists only when
\[
  \det(A-\lambda I)=0.
\]

\begin{example}
Let
\[
  A=\begin{pmatrix}2&1\\1&2\end{pmatrix}.
\]
Then
\[
  \det(A-\lambda I)=(2-\lambda)^2-1.
\]
The eigenvalues are \(3\) and \(1\).  For \(\lambda=3\), an eigenvector is \((1,1)\); for \(\lambda=1\), an eigenvector is \((1,-1)\).  The matrix stretches the symmetric direction by \(3\) and the antisymmetric direction by \(1\).
\end{example}

\section{Complex vector spaces and Hilbert spaces}

Quantum states require complex numbers.  For complex vectors, the inner product is
\[
  \braket{u}{v}=\sum_i u_i^*v_i.
\]
The conjugate on the first slot ensures
\[
  \braket{v}{v}=\sum_i |v_i|^2\ge0.
\]
A Hilbert space is a vector space with an inner product, completed so that limits of convergent sequences remain in the space.  In this book we will mostly use finite-dimensional intuition and the function-space example
\[
  \braket{f}{g}=\int f^*(x)g(x)\dd x.
\]

The adjoint \(A^\dagger\) is defined by
\[
  \braket{u}{Av}=\braket{A^\dagger u}{v}.
\]
In matrix form, \(A^\dagger\) is transpose plus complex conjugation:
\[
  (A^\dagger)_{ij}=A_{ji}^*.
\]
If \(A=A^\dagger\), then \(A\) is Hermitian.  Hermitian matrices have real eigenvalues:
\[
  A\ket{v}=\lambda\ket{v}
  \quad\Rightarrow\quad
  \bra{v}A\ket{v}=\lambda\braket{v}{v}.
\]
But also
\[
  \bra{v}A\ket{v}=\braket{Av}{v}=\lambda^*\braket{v}{v},
\]
so \(\lambda=\lambda^*\) when \(\ket v\ne0\).  Hermitian operators become observables.

\section{Dirac notation}

A ket \(\ket\psi\) is a vector.  A bra \(\bra\psi\) is the corresponding linear functional obtained by conjugate transpose.  The inner product is \(\braket{\phi}{\psi}\).  The outer product
\[
  \ket\phi\bra\psi
\]
is an operator, because it acts on \(\ket\chi\) as
\[
  \ket\phi\bra\psi\ket\chi
  =
  \ket\phi\,\braket{\psi}{\chi}.
\]
If \(\{\ket n\}\) is an orthonormal basis, then
\[
  \braket{m}{n}=\delta_{mn},
  \qquad
  \sum_n\ket n\bra n=\id.
\]
The second equation says that every vector can be rebuilt from its components:
\[
  \ket\psi=\sum_n\ket n\braket{n}{\psi}.
\]

\section{Operators on functions}

The derivative is a linear operator:
\[
  D(\alpha f+\beta g)=\alpha Df+\beta Dg.
\]
With boundary terms vanishing,
\[
  \int f^*(x)\frac{\dd g}{\dd x}\dd x
  =
  -\int \left(\frac{\dd f}{\dd x}\right)^*g(x)\dd x.
\]
Thus \(D^\dagger=-D\), and \(-\ii D\) is Hermitian:
\[
  (-\ii D)^\dagger=\ii D^\dagger=-\ii D.
\]
This is why the momentum operator in one dimension is
\[
  \hat p=-\ii\hbar\frac{\dd}{\dd x}.
\]
It has plane-wave eigenfunctions:
\[
  \hat p e^{\ii kx}=\hbar k e^{\ii kx}.
\]

\section*{Exercises}
\addcontentsline{toc}{section}{Exercises}

\begin{exercise}
For
\[
  B=\begin{pmatrix}0&-\ii\\ \ii&0\end{pmatrix},
\]
show that \(B\) is Hermitian and find its eigenvalues.
\begin{answer}
\(B^\dagger=B\).  The characteristic polynomial is \(\lambda^2-1\), so \(\lambda=\pm1\).
\end{answer}
\end{exercise}

\begin{exercise}
Prove that if \(A\) and \(B\) are Hermitian, then \(\ii[A,B]\) is Hermitian.
\begin{answer}
\([A,B]^\dagger=(AB-BA)^\dagger=BA-AB=-[A,B]\).  Therefore \((\ii[A,B])^\dagger=-\ii(-[A,B])=\ii[A,B]\).
\end{answer}
\end{exercise}
